\section{Related Works}
\label{sec:related_works}

% Narasi pengantar untuk tabel
Recent studies have extensively explored the application of machine learning in financial fraud detection, focusing heavily on adapting to shifting transaction patterns known as concept drift. Furthermore, addressing the extreme class imbalance inherent in real-world fraud scenarios remains a formidable challenge. Table \ref{tab:research_gap} summarizes the state-of-the-art literature and highlights the specific research gaps filled by our proposed method.

% Tabel Research Gap
\begin{table}[htbp]
    \centering
    \caption{Comparison of State-of-the-Art Literature and Our Proposed Method}
    \label{tab:research_gap}
    \resizebox{\textwidth}{!}{%
    \begin{tabular}{|l|l|c|c|p{6cm}|}
        \hline
        \textbf{Reference / Year} & \textbf{Proposed Method} & \textbf{Drift Adaptation?} & \textbf{Imbalance Handling?} & \textbf{Limitations / Research Gap} \\ \hline
        
        Author A et al. (2025) \cite{ref_a} & Ensemble-Based Fraud Detection (Deep Auto-Encoder / SVM) & \ding{55} No & \ding{51} Yes & Model is trained offline (batch processing), making it highly susceptible to obsolescence when new fraud patterns emerge. \\ \hline
        
        Author B et al. (2024) \cite{ref_b} & Real-Time Fraud Detection using Kafka \& ML Ensemble & \ding{51} Yes & \ding{55} No & Focuses on infrastructure latency but employs standard ensemble models that suffer from severe performance degradation on the minority class (<0.2\%). \\ \hline
        
        Author C et al. (2025) \cite{ref_c} & Incremental Gradient Boosting Trees with Adaptive Windowing & \ding{51} Yes & \ding{55} No & Excellent at detecting concept drift, but exhibits low fraud recall in real-world distributions due to the lack of cost-sensitive objective functions. \\ \hline
        
        Author D et al. (2026) \cite{ref_d} & Two-stage cost-sensitive learning for data streams & \ding{51} Yes & \ding{51} Yes & Utilizes a two-stage preprocessing architecture that incurs significant computational overhead, rendering it impractical for high-velocity real-time payment streams. \\ \hline
        
        \textbf{Proposed Method} & \textbf{Cost-Sensitive Adaptive Random Forest (CS-ARF)} & \textbf{\ding{51} Yes} & \textbf{\ding{51} Yes} & \textbf{Eliminates two-stage computational overhead by utilizing Implicit Online Oversampling directly within the ARF algorithm, resolving both drift and imbalance simultaneously in a single pass.} \\ \hline
    \end{tabular}%
    }
\end{table}

% Narasi penjelasan celah (gap) penelitian
As demonstrated in Table \ref{tab:research_gap}, although many recent studies attempt to tackle concept drift using adaptive windowing mechanisms \cite{ref_c}, they predominantly overlook the extreme class imbalance where fraud constitutes less than 0.2\% of transactions. Conversely, methods that attempt to solve both problems simultaneously \cite{ref_d} often rely on two-stage preprocessing architectures, which are computationally expensive for real-time processing. 

To bridge this crucial gap, our research proposes the direct integration of Cost-Sensitive Learning—via Implicit Online Oversampling—into an Adaptive Random Forest architecture. This approach enables the model to swiftly adapt to novel fraud patterns while imposing a massive algorithmic penalty against misclassifying the minority class, all executed within a highly efficient single-pass streaming computation.
